\documentclass[11pt]{article}

\usepackage[a4paper,margin=1in]{geometry}
\usepackage{amsmath,amssymb,amsthm}
\usepackage{booktabs}
\usepackage{array}
\usepackage{longtable}
\usepackage{hyperref}

\newcommand{\F}{\mathbb F}
\newcommand{\PG}{\mathrm{PG}}
\newcommand{\RL}{\mathrm{RL}}
\newcommand{\wt}{\mathrm{wt}}

\theoremstyle{plain}
\newtheorem{theorem}{Theorem}
\newtheorem{proposition}{Proposition}
\newtheorem{corollary}{Corollary}

\theoremstyle{definition}
\newtheorem{definition}{Definition}
\newtheorem{remark}{Remark}

\title{Covering radii and deep-hole syndromes of\\
length-eight Roth--Lempel codes}
\author{Yong Zhang\\
School of Mathematics and Statistics, Yancheng Normal University\\
Yancheng, Jiangsu 224002, China
}
\date{}

\begin{document}
\maketitle

\begin{abstract}
We determine covering radii and deep-hole syndromes for length-eight,
dimension-four Roth--Lempel codes.  They have the form
\(C=\RL_{4,\delta}(S)\), where \(|S|=6\).
First, \(C\) is always MDS or NMDS, and it is NMDS precisely when
\(\delta\) is the sum of three distinct elements of \(S\).  Our main result
is the sharp threshold \(q\ge16\Rightarrow \rho(C)=4\).
For the remaining admissible fields, the radius-three codes are classified
completely up to affine equivalence: their numbers are respectively
\(2,3,14,32,7\) over \(\F_7,\F_8,\F_9,\F_{11},\F_{13}\), and no other
radius-three case exists.  The proof translates \(\rho(C)\le3\) into a
finite-geometric covering condition: the \(56\) planes spanned by triples of
parity-check columns must cover \(\PG(3,q)\).  An uncovered point is both a
deep-hole syndrome and a projection center producing a plane \(8\)-arc.
An explicit twisted-cubic construction proves the threshold theoretically
for every \(q\ge23\), and for every NMDS member already when \(q\ge19\);
only thirty affine representatives over \(q=16,17,19\) require explicit
finite certificates.  We further give a \(56\)-factor polynomial whose
nonvanishing characterizes all deep-hole syndromes when \(q\ge16\), an
explicit distance stratification in every admissible field, and closed
formulas for all projective syndrome shells, including
\(m_2(C)=28(q-1)-3r_S(\delta)\).
\end{abstract}

\section{Introduction}

Covering radii and deep holes of linear codes are classical questions in
coding theory \cite{Cohen1997,Huffman2003,MWS1977}.  For MDS and NMDS codes,
their syndrome geometry is naturally expressed by point configurations and
arcs in finite projective spaces \cite{Ball2015,DodunekovLandgev1995}.
Roth--Lempel codes are a classical source of non-Reed--Solomon MDS codes
\cite{RothLempel1989}.

The recent literature closest to the present problem has developed in two
directions.  Han and Fan, and subsequently Xu and Zhou, investigate NMDS
conditions and weight distributions for Roth--Lempel codes
\cite{HanFan2023,XuZhou2026}; Liang and Liao construct further NMDS classes
from the Roth--Lempel matrix \cite{LiangLiao2026}.  In parallel, covering
radii and deep holes have become structural tools for extended twisted-GRS
and non-GRS MDS--NMDS constructions
\cite{LiZhuSun2025,LiLuLingLam2026}.  Distinct from these works, this paper
determines the complete covering-radius threshold and explicit deep-hole
syndrome criterion for the fixed length-eight, dimension-four
Roth--Lempel family.

The problem is particularly sharp in this family.  Since every
\(\RL_{4,\delta}(S)\) considered here is an \([8,4]_q\) code, the redundancy
bound gives \(\rho(C)\le4\).  Thus the only question is whether the covering
radius drops from \(4\) to \(3\).  We answer this question for every prime
power \(q\ge7\), and at the same time determine the corresponding deep-hole
syndromes.

The main results are as follows.
\begin{enumerate}
\item Every \(\RL_{4,\delta}(S)\), \(|S|=6\), is MDS or NMDS.  Moreover,
it is NMDS if and only if
\(\delta=a+b+c\) for three distinct elements \(a,b,c\in S\)
(Proposition~\ref{prop:exhaustion}).
\item The covering-radius threshold is exact:
\[
 \rho\bigl(\RL_{4,\delta}(S)\bigr)=4
 \qquad\text{for every prime power }q\ge16.
\]
Below the threshold, all radius-three affine classes are determined:
there are \(2,3,14,32,7\) over \(q=7,8,9,11,13\), respectively
(Theorems~\ref{thm:complete} and~\ref{thm:threshold}).
\item A syndrome has distance at most \(3\) exactly when its projective point
lies in one of the \(56\) planes spanned by triples of parity-check columns
(Proposition~\ref{prop:criterion}).  Consequently, uncovered points are
simultaneously deep-hole syndromes and projection centers that produce plane
\(8\)-arcs (Proposition~\ref{prop:projection}).
\item For \(q\ge16\), all deep-hole syndromes are characterized by the
nonvanishing of one explicit product of \(56\) linear factors
(Proposition~\ref{prop:explicit-product}).  The twisted-cubic subfamily of
centers already proves existence for all \(q\ge23\), and for all NMDS codes
when \(q\ge19\) (Proposition~\ref{prop:twisted-centers}).
\item The complete syndrome distance stratification is explicit in every
admissible field, and the projective shell sizes admit closed formulas.  In
particular, the second shell is
\[
 m_2(C)=28(q-1)-3r_S(\delta),
\]
where \(r_S(\delta)\) counts the three-subsets of \(S\) with sum \(\delta\)
(Theorem~\ref{thm:distance-stratification} and
Corollary~\ref{cor:closed-shells}).
\end{enumerate}

These conclusions concern the existence and explicit description of deep
holes in a classical non-GRS MDS/NMDS family; they do not require introducing
a new code definition.  They nevertheless provide an infinite supply of
fully certified \([8,4]_q\) Roth--Lempel seed codes with maximal covering
radius for deep-hole-based extension constructions.  The finite part of the
proof uses exact Magma enumeration, not randomized search.  The proof-critical
computations, certificates, and independent cross-checks are separated and
documented in Section~\ref{sec:reproducibility}.

\section{Roth--Lempel codes}

Let \(3\leq k\leq n\leq q\), let \(S=\{a_1,\ldots,a_n\}\subseteq \F_q\), and
let \(\delta\in\F_q\).  The Roth--Lempel code
\(\RL_{k,\delta}(S)\) is the \([n+2,k]_q\) code
\[
  \bigl\{
    (f(a_1),\ldots,f(a_n),f_{k-1},f_{k-2}+\delta f_{k-1}):
    f(x)=\sum_{i=0}^{k-1} f_i x^i\in\F_q[x]
  \bigr\}.
\]
Equivalently, it has generator matrix
\[
  G_{\RL_{k,\delta}(S)}=
  \begin{pmatrix}
    1 & 1 & \cdots & 1 & 0 & 0\\
    a_1 & a_2 & \cdots & a_n & 0 & 0\\
    \vdots & \vdots & \ddots & \vdots & \vdots & \vdots\\
    a_1^{k-3} & a_2^{k-3} & \cdots & a_n^{k-3} & 0 & 0\\
    a_1^{k-2} & a_2^{k-2} & \cdots & a_n^{k-2} & 0 & 1\\
    a_1^{k-1} & a_2^{k-1} & \cdots & a_n^{k-1} & 1 & \delta
  \end{pmatrix}.
\]
We focus on \(k=4\) and \(n=6\), so that the resulting code has parameters
\([8,4]_q\).
We use the standard convention that an \([N,K]_q\) code is NMDS when
\(d(C)=N-K\) and \(d(C^\perp)=K\), following
Dodunekov and Landgev \cite{DodunekovLandgev1995}.

The natural affine change of variables
\[
  x\mapsto ax+b,\qquad a\in\F_q^\ast,\ b\in\F_q,
\]
acts on parameters by
\[
  S\mapsto aS+b,\qquad \delta\mapsto a\delta+3b.
\]
All classifications below are taken modulo this affine action.

\begin{proposition}[MDS/NMDS exhaustion for the family]\label{prop:exhaustion}
For every prime power \(q\ge7\), every code
\(\RL_{4,\delta}(S)\), with \(S\subseteq\F_q\) and \(|S|=6\), is either MDS
or NMDS.  More precisely, it is NMDS if and only if
\[
 \delta=a+b+c\quad\hbox{for some distinct }a,b,c\in S.
\]
\end{proposition}

\begin{proof}
Write the last two generator columns as
\(e=(0,0,0,1)^{\mathsf T}\) and
\(f=(0,0,1,\delta)^{\mathsf T}\), and the six evaluation columns as
\(v(a)=(1,a,a^2,a^3)^{\mathsf T}\), \(a\in S\).  Every three columns are
independent: this is the Vandermonde determinant for three evaluation
columns, and with one or two of \(e,f\) it follows respectively from the
rank-two first-coordinate projection of two distinct evaluation columns, or
from the nonzero first coordinate of a single evaluation column together
with the independence of \(e,f\).  Every
five columns have rank four.  This is immediate if at least four evaluation
columns occur; if exactly three occur, then both \(e,f\) occur and quotienting
by their span leaves the rank-two matrix with columns \((1,a)^{\mathsf T}\)
for three distinct \(a\)'s.  Thus either every four columns are independent,
giving an MDS code, or some four are dependent while every three are
independent and every five have rank four, giving an NMDS code.
The only four-column sets that can become dependent are
\(\{v(a),v(b),v(c),f\}\): their determinant is a nonzero Vandermonde
factor times \(\delta-(a+b+c)\).  This proves the final assertion.
\end{proof}

\begin{proposition}[Explicit duality]\label{prop:duality}
Put \(\sigma=\sum_{a\in S}a\), \(\theta=\sigma-\delta\), and
\[
  \lambda_i=\left(\prod_{j\ne i}(a_i-a_j)\right)^{-1}.
\]
For \(C=\RL_{4,\delta}(S)\), a parity-check matrix is
\[
H=
\begin{pmatrix}
\lambda_1&\cdots&\lambda_6&0&0\\
\lambda_1a_1&\cdots&\lambda_6a_6&0&0\\
\lambda_1a_1^2&\cdots&\lambda_6a_6^2&-1&0\\
\lambda_1a_1^3&\cdots&\lambda_6a_6^3&\delta-\sigma&-1
\end{pmatrix}.
\]
Consequently \(C^\perp\) is monomially equivalent to
\(\RL_{4,\theta}(S)\).
\end{proposition}

\begin{proof}
The Lagrange denominator identities give
\[
 \sum_{i=1}^6\lambda_i a_i^r=0\quad(0\le r\le4),
 \qquad
 \sum_{i=1}^6\lambda_i a_i^5=1,
 \qquad
 \sum_{i=1}^6\lambda_i a_i^6=\sigma.
\]
A direct multiplication by \(G_{\RL_{4,\delta}(S)}\) therefore gives
\(GH^{\mathsf T}=0\), and \(H\) has rank \(4\).  Scaling its first six
columns by \(\lambda_i^{-1}\), scaling the last two by \(-1\), and swapping
the last two columns produces the eight columns
\[
 (1,a,a^2,a^3)^{\mathsf T}\ (a\in S),\qquad
 (0,0,0,1)^{\mathsf T},\qquad(0,0,1,\theta)^{\mathsf T},
\]
which are precisely the generator columns of \(\RL_{4,\theta}(S)\).
\end{proof}

\begin{corollary}[Dual compatibility of the MDS criterion]
\label{cor:duality-mds}
Let \(r_S(\delta)\) denote the number of three-subsets \(T\subset S\) with
\(\sum_{a\in T}a=\delta\).  Then \(C\) is MDS if and only if
\(r_S(\delta)=0\), and it is NMDS if and only if \(r_S(\delta)>0\).
Moreover,
\[
  r_S(\delta)=r_S\left(\sum_{a\in S}a-\delta\right).
\]
\end{corollary}

\begin{proof}
The type criterion is Proposition~\ref{prop:exhaustion}.  Complementation
maps a three-subset of sum \(\delta\) bijectively to a three-subset of sum
\(\sum_{a\in S}a-\delta\), in agreement with
Proposition~\ref{prop:duality}.
\end{proof}

\begin{proposition}[The second syndrome shell]\label{prop:second-shell}
For \(C=\RL_{4,\delta}(S)\), let \(r_S(\delta)\) be as in
Corollary~\ref{cor:duality-mds}.  The number \(m_2(C)\) of projective
syndrome points at distance exactly \(2\) satisfies
\[
  m_2(C)=28(q-1)-3r_S(\delta).
\]
\end{proposition}

\begin{proof}
The eight parity-check points determine \(28\) secant lines.  Since every
three columns are independent, two secants sharing a parity-check point have
no further point in common.  If no four columns are dependent, the union of
the secants has
\[
  8+28(q-1)
\]
points, and hence \(m_2(C)=28(q-1)\).

By Proposition~\ref{prop:duality}, every dependent four-set of parity-check
points has the form
\[
  \{V(a),V(b),V(c),E\},\qquad a+b+c=\theta.
\]
Equivalently, it is indexed by one of the
\(r_S(\delta)=r_S(\theta)\) complementary three-subsets.
The three pairs of opposite secants in such a plane meet in three
non-column points, decreasing the second shell by \(3\).  Contributions from
two different dependent four-sets cannot coincide.  Each diagonal point
lies on one secant through \(E\).  If it came from two circuits, this
secant would be the same line \(\langle E,V(a)\rangle\), and the opposite
secants would meet at the same non-column point.  If they share an endpoint,
their intersection is a column; if they are disjoint, their four endpoints
are coplanar, contrary to the Vandermonde determinant.
Therefore the total decrement is exactly \(3r_S(\delta)\).
\end{proof}

\section{A trisecant-plane criterion}

Let \(C\) be an \([8,4]_q\) linear code and let
\[
  H_1,\ldots,H_8\in \PG(3,q)
\]
be the projective columns of a parity-check matrix \(H\) of \(C\).  A
syndrome \(s\in\F_q^4\) has a coset representative of weight at most \(t\) if
and only if its projective point lies in the projective span of at most \(t\)
columns of \(H\).  For \(t=3\) this gives the following criterion.

\begin{proposition}[Trisecant-plane criterion]\label{prop:criterion}
Let \(C\) be an \([8,4]_q\) linear code with parity-check columns
\(H_1,\ldots,H_8\in\PG(3,q)\).  Then
\[
  \rho(C)\le 3
  \quad\Longleftrightarrow\quad
  \PG(3,q)=
  \bigcup_{1\le i<j<\ell\le 8}
  \langle H_i,H_j,H_\ell\rangle .
\]
Since \(\rho(C)\le 4\) for every \([8,4]_q\) linear code, failure of this
cover is equivalent to \(\rho(C)=4\).
\end{proposition}

\begin{proof}
The syndrome of a word supported on a set \(I\subseteq\{1,\ldots,8\}\) is a
linear combination of the columns \(\{H_i:i\in I\}\).  Thus a projective
syndrome has a representative of weight at most \(3\) precisely when it lies
in the span of one, two, or three parity-check columns.  These smaller spans
are contained in the spans of triples.  Hence every syndrome has a
representative of weight at most \(3\) if and only if the \(56\) triple spans
cover \(\PG(3,q)\).  The final assertion is the standard redundancy bound
\cite{Cohen1997,Huffman2003}.
\end{proof}

\begin{proposition}[Projection centers]\label{prop:projection}
With the notation of Proposition~\ref{prop:criterion}, a projective point
\(P\in\PG(3,q)\) is not covered by the \(56\) triple spans if and only if
projection from \(P\) maps \(H_1,\ldots,H_8\) to eight points of a plane with
no coincident pair and no collinear triple.  In other words, uncovered
syndrome points are exactly projection centers producing plane \(8\)-arcs.
\end{proposition}

\begin{proof}
Two projected columns coincide exactly when \(P\) lies on the line joining the
corresponding two columns.  Three projected columns are collinear exactly when
\(P\) lies in the plane spanned by the corresponding three columns.  Therefore
absence from all triple spans is equivalent to the projected configuration
having neither a collision nor a collinear triple.
\end{proof}

Thus the uncovered syndrome points are not only deep-hole cosets; they are
also constructive centers for producing plane \(8\)-arcs from the
Roth--Lempel parity-check configuration.  The explicit centers constructed
below therefore give a finite-geometric witness for the existence of such
projected arcs in the non-exceptional range.

\begin{proposition}[Explicit uncovered-syndrome criterion]
\label{prop:explicit-product}
Let \(C=\RL_{4,\delta}(S)\), let
\(\theta=\sum_{a\in S}a-\delta\), and write a projective syndrome relative
to the parity-check matrix in Proposition~\ref{prop:duality} as
\(P=(p_0,p_1,p_2,p_3)\).  Define
\[
\begin{aligned}
A_a(P)&=p_1-ap_0,\\
B_{ab}(P)&=p_2-(a+b)p_1+abp_0,\\
C_{ab}(P)&=p_3-\theta p_2+
 \{ab+(a+b)(\theta-a-b)\}p_1-ab(\theta-a-b)p_0,\\
D_{abc}(P)&=p_3-(a+b+c)p_2+(ab+ac+bc)p_1-abcp_0
\end{aligned}
\]
and
\[
 \Phi_{S,\delta}(P)=
 \prod_{a\in S}A_a(P)
 \prod_{\{a,b\}\subset S}B_{ab}(P)C_{ab}(P)
 \prod_{\{a,b,c\}\subset S}D_{abc}(P).
\]
Then \(P\) is outside every triple span of parity-check columns if and only
if \(\Phi_{S,\delta}(P)\ne0\).  Equivalently, a syndrome has coset distance
\(4\) if and only if its projective point satisfies this nonvanishing
condition.
\end{proposition}

\begin{proof}
By Proposition~\ref{prop:duality}, column scaling does not affect the
projective spans, and the parity-check points may be written as
\[
 V(a)=(1,a,a^2,a^3)^{\mathsf T}\quad(a\in S),\qquad
 E=(0,0,1,\theta)^{\mathsf T},\qquad F=(0,0,0,1)^{\mathsf T}.
\]
The \(56\) triple spans split into four types.  Their plane equations at
\(P\), up to nonzero scalar multiples, are respectively
\[
\begin{array}{c|c|c}
\hbox{triple type}&\hbox{number}&\hbox{plane equation}\\ \hline
\{V(a),E,F\}&6&A_a(P)=0\\
\{V(a),V(b),F\}&15&B_{ab}(P)=0\\
\{V(a),V(b),E\}&15&C_{ab}(P)=0\\
\{V(a),V(b),V(c)\}&20&D_{abc}(P)=0.
\end{array}
\]
For the third row, the cubic polynomial defining the plane has roots
\(a,b,\theta-a-b\).  Hence \(P\) avoids all triple planes precisely when
none of the displayed \(56\) factors vanishes.  The final statement follows
from Proposition~\ref{prop:criterion}.
\end{proof}

\begin{proposition}[Twisted-cubic deep-hole centers]
\label{prop:twisted-centers}
Let \(C=\RL_{4,\delta}(S)\), put
\(\theta=\sum_{a\in S}a-\delta\), and set
\[
 {\cal B}_{S,\delta}
 =S\cup\{\theta-a-b:\{a,b\}\subset S\}.
\]
For \(x\in\F_q\), the point \(V(x)=(1,x,x^2,x^3)\) has coset distance
\(4\) if and only if \(x\notin{\cal B}_{S,\delta}\).  Let \(m_4(C)\)
denote the number of projective syndromes at distance \(4\), and let
\(r=r_S(\delta)\).  Then
\[
  \lvert{\cal B}_{S,\delta}\rvert\le21-3r,\qquad
  m_4(C)\ge q-21+3r.
\]
Consequently:
\[
\begin{array}{ll}
q\ge23&\Longrightarrow\quad \rho(C)=4,\quad
 m_4(C)\ge q-21;\\
C\hbox{ NMDS and }q\ge19&\Longrightarrow\quad \rho(C)=4,\quad
 m_4(C)\ge q-18;\\
r\ge2\hbox{ and }q\ge16&\Longrightarrow\quad \rho(C)=4,\quad
 m_4(C)\ge q-15;\\
r\ge3\hbox{ and }q\ge13&\Longrightarrow\quad \rho(C)=4,\quad
 m_4(C)\ge q-12.
\end{array}
\]
The corresponding lower bound for the number of deep-hole words is
\(q^4(q-1)(q-21+3r)\) whenever the displayed lower bound is positive.
\end{proposition}

\begin{proof}
Substitution of \(P=V(x)\) into the factors of
Proposition~\ref{prop:explicit-product} gives
\[
\begin{aligned}
A_a(V(x))&=x-a,\\
B_{ab}(V(x))&=(x-a)(x-b),\\
C_{ab}(V(x))&=(x-a)(x-b)(x-\theta+a+b),\\
D_{abc}(V(x))&=(x-a)(x-b)(x-c).
\end{aligned}
\]
Hence \(\Phi_{S,\delta}(V(x))\ne0\) precisely when
\(x\notin{\cal B}_{S,\delta}\).

By Corollary~\ref{cor:duality-mds}, exactly \(r\) three-subsets of \(S\)
have sum \(\theta\).  A pair \(\{a,b\}\subset S\) satisfies
\(\theta-a-b\in S\) precisely when it lies in one of these three-subsets;
each such subset supplies three pairs, and each qualifying pair has a
unique third member.  Thus exactly \(3r\) of the \(15\) indexed pair values
already belong to \(S\).  Any additional coincidence only decreases the
forbidden set, giving
\[
 \lvert{\cal B}_{S,\delta}\rvert\le 6+(15-3r)=21-3r.
\]
There are therefore at least \(q-21+3r\) twisted-cubic centers whenever
this number is positive.  The assertions follow from
Proposition~\ref{prop:explicit-product} and
Proposition~\ref{prop:deep-census}.
\end{proof}

\begin{proposition}[Affine covariance of the explicit centers]
\label{prop:affine-covariance}
Let \(S'=uS+v\) and \(\delta'=u\delta+3v\), where
\(u\in\F_q^\ast\) and \(v\in\F_q\).  If
\(\theta=\sum_{a\in S}a-\delta\) and
\(\theta'=\sum_{a'\in S'}a'-\delta'\), then
\[
 \theta'=u\theta+3v,\qquad
 {\cal B}_{S',\delta'}=u{\cal B}_{S,\delta}+v.
\]
In particular, the number of twisted-cubic centers supplied by
Proposition~\ref{prop:twisted-centers} is invariant under the affine
classification action.
\end{proposition}

\begin{proof}
Since \(|S|=6\),
\[
 \theta'=u\sum_{a\in S}a+6v-u\delta-3v=u\theta+3v.
\]
Furthermore,
\[
 \theta'-(ua+v)-(ub+v)=u(\theta-a-b)+v,
\]
and the asserted transformation of \({\cal B}_{S,\delta}\) follows.
\end{proof}

\begin{remark}[Auxiliary boundary families]\label{rem:boundary-families}
For finite boundary certificates it is also useful to consider
\[
 W(x,\mu)=(1,x,x^2,x^3+\mu),\qquad Q(t,u)=(0,1,t,u).
\]
Substitution in the factors of Proposition~\ref{prop:explicit-product}
shows that, once the first parameter avoids the corresponding secant
obstructions, each remaining triple-plane obstruction excludes one value
of the second parameter.  These families compress the boundary
certificates below, but their elementary count does not improve the
uniform threshold in Proposition~\ref{prop:twisted-centers}.
\end{remark}

\begin{theorem}[Explicit syndrome distance stratification]
\label{thm:distance-stratification}
Continue with the notation of Proposition~\ref{prop:explicit-product}, and
let \(P\ne0\) be a syndrome point.  Write \(\mathcal L_2(P)\) for the
following disjunction:
\[
\begin{array}{ll}
p_0=p_1=0; & \\[1mm]
p_1-ap_0=p_2-a^2p_0=0
  &\hbox{for some }a\in S;\\[1mm]
p_1-ap_0=0,\quad
p_3-\theta p_2-(a^3-\theta a^2)p_0=0
  &\hbox{for some }a\in S;\\[1mm]
B_{ab}(P)=0,\quad
p_3-(a+b)p_2+abp_1=0
  &\hbox{for some }\{a,b\}\subset S.
\end{array}
\]
Let \(\mathcal X=\{V(a):a\in S\}\cup\{E,F\}\), with \(V(a),E,F\) as in
the proof of Proposition~\ref{prop:explicit-product}.  Then the coset
distance represented by \(P\) is exactly
\[
\begin{array}{c|c}
\hbox{distance}&\hbox{condition}\\ \hline
1&P\in\mathcal X,\\
2&P\notin\mathcal X\ \hbox{and}\ \mathcal L_2(P),\\
3&\neg\mathcal L_2(P)\ \hbox{and}\ \Phi_{S,\delta}(P)=0,\\
4&\Phi_{S,\delta}(P)\ne0.
\end{array}
\]
In particular, this gives an explicit deep-hole test also for every
radius-three exceptional class in Theorem~\ref{thm:complete}: its deep-hole
syndromes are precisely the points satisfying \(\neg\mathcal L_2(P)\).
\end{theorem}

\begin{proof}
The four lines in the definition of \(\mathcal L_2(P)\) are respectively
the equations for the \(28\) secant lines
\[
 \langle E,F\rangle,\quad
 \langle V(a),F\rangle,\quad
 \langle V(a),E\rangle,\quad
 \langle V(a),V(b)\rangle.
\]
Thus \(\mathcal L_2(P)\) is equivalent to syndrome distance at most \(2\).
Proposition~\ref{prop:explicit-product} says that
\(\Phi_{S,\delta}(P)=0\) is equivalent to distance at most \(3\), while
nonvanishing is equivalent to distance \(4\).  Subtracting successive
layers proves the table.  Finally, the trisecant cover is complete in a
radius-three class, so only the first three layers remain.
\end{proof}

\begin{proposition}[Projective shells and deep-hole census]
\label{prop:deep-census}
Let \(C\) be an \([8,4]_q\) code and, for \(1\le t\le4\), let \(m_t(C)\)
denote the number of projective syndrome points whose minimum expression as a
span of parity-check columns uses exactly \(t\) columns.  Then
\[
  \rho(C)=\max\{t:m_t(C)>0\},
  \qquad
  |\mathrm{DH}(C)|=q^4(q-1)m_{\rho(C)}(C),
\]
where \(\mathrm{DH}(C)\) is the set of deep-hole words of \(C\).
In particular, if \(\rho(C)=4\), then \(m_4(C)\) is exactly the number of
uncovered points in Proposition~\ref{prop:criterion}.
\end{proposition}

\begin{proof}
A nonzero syndrome has coset distance \(t\) precisely when its projective
point belongs to the \(t\)-th shell.  Each projective point accounts for
\(q-1\) nonzero scalar syndromes, and each syndrome coset contains
\(|C|=q^4\) words.  The final assertion follows from the trisecant-plane
criterion.
\end{proof}

\begin{corollary}[Closed shell formula for the family]
\label{cor:closed-shells}
Let \(C=\RL_{4,\delta}(S)\), set \(r=r_S(\delta)\), and write
\(h=m_4(C)\) for the number of uncovered projective syndrome points.  Then
\[
 \bigl(m_1(C),m_2(C),m_3(C),m_4(C)\bigr)
 =
 \bigl(8,\ 28(q-1)-3r,\ q^3+q^2-27q+21+3r-h,\ h\bigr).
\]
In particular, if \(\rho(C)=3\), then
\[
 |\mathrm{DH}(C)|
 =q^4(q-1)\bigl(q^3+q^2-27q+21+3r_S(\delta)\bigr).
\]
\end{corollary}

\begin{proof}
The eight projective parity-check columns are distinct, so \(m_1(C)=8\).
Proposition~\ref{prop:second-shell} gives \(m_2(C)\), and
\[
 m_1(C)+m_2(C)+m_3(C)+m_4(C)=q^3+q^2+q+1.
\]
Solving for \(m_3(C)\) gives the stated vector.  If \(\rho(C)=3\), then
\(h=0\), and Proposition~\ref{prop:deep-census} converts \(m_3(C)\) to
the number of deep-hole words.
\end{proof}

\begin{corollary}[An infinite large-field class]\label{cor:large-q}
Let \(C\) be any \([8,4]_q\) linear code.  If \(q\ge 56\), then
\(\rho(C)=4\).  In particular, for every prime power \(q\ge 56\), every
Roth--Lempel code \(\RL_{4,\delta}(S)\) with \(|S|=6\) has covering radius
\(4\).
\end{corollary}

\begin{proof}
A plane in \(\PG(3,q)\) contains \(q^2+q+1\) points, while
\(\PG(3,q)\) contains \(q^3+q^2+q+1\) points.  The \(56\) planes spanned by
triples of parity-check columns cover at most \(56(q^2+q+1)\) points.  For
\(q\ge 56\),
\[
  q^3+q^2+q+1>56(q^2+q+1).
\]
Thus these planes cannot cover \(\PG(3,q)\), so
Proposition~\ref{prop:criterion} gives \(\rho(C)>3\).  Since an \([8,4]_q\)
code has redundancy \(4\), the redundancy bound gives \(\rho(C)\le 4\).
Therefore \(\rho(C)=4\).
\end{proof}

\begin{proposition}[MDS and NMDS refinements]\label{prop:large-q-refined}
Let \(C\) be an \([8,4]_q\) code.  If \(C\) is MDS and \(q\ge54\), then
\(\rho(C)=4\).  If \(C\) is NMDS and \(q\ge53\), then
\(\rho(C)=4\).
\end{proposition}

\begin{proof}
Assume first that \(C\) is MDS.  Its \(56\) triple planes are distinct.  At
each of the eight columns, \(21\) triple planes meet, contributing at least
\(20\) repeated incidences.  Further, each of the \(28\) secant lines is
contained in the six triple planes obtained by adjoining one of the remaining
columns.  The \(28\) sets of \(q-1\) non-column points on these secants are
pairwise disjoint: an intersection would put four columns in a plane.  Hence
the union of the triple planes has at most
\[
  56(q^2+q+1)-8\cdot20-28(q-1)\cdot5
  =56q^2-84q+36
\]
points.  This is smaller than \(q^3+q^2+q+1\) for \(q\ge54\).

If \(C\) is NMDS, some four columns are coplanar while every triple is
independent.  The four triple spans arising from that four-set coincide, so
there are at most \(53\) distinct triple planes.  Since
\[
  53(q^2+q+1)<q^3+q^2+q+1\qquad(q\ge53),
\]
they do not cover \(\PG(3,q)\).  In either case,
Proposition~\ref{prop:criterion} and the redundancy bound yield
\(\rho(C)=4\).
\end{proof}

\section{Complete affine classifications over small fields}

For each prime power
\[
q\in\{7,8,9,11,13,16,17,19,23,25,27,29,31,32,37,41,43,47,49,53\}
\]
we enumerated all pairs \((S,\delta)\) with
\(|S|=6\), reduced them under the affine action above, computed the covering
radius or certified radius \(4\) by an explicit uncovered syndrome point for
\(\RL_{4,\delta}(S)\).  This normalization is the natural affine
specialization of the finite-geometric equivalence used for point
configurations and arcs in projective spaces \cite{Ball2015}.

\begin{theorem}[Complete small-field classification]\label{thm:complete}
For the Roth--Lempel codes \(\RL_{4,\delta}(S)\) with \(|S|=6\), the complete
affine classifications over every prime-power field \(\F_q\) with
\(7\le q\le53\) are as follows.  No class in the table has radius below \(3\).
{\scriptsize
\[
\begin{array}{c|r|r|r|r|r|r}
q & \hbox{raw cases} & \hbox{affine classes}
  & \rho=3\hbox{ classes} & \rho=3\hbox{ cases}
  & \rho=4\hbox{ classes} & \rho=4\hbox{ cases}\\
\hline
7 & 49 & 2 & 2 & 49 & 0 & 0\\
8 & 224 & 4 & 3 & 168 & 1 & 56\\
9 & 756 & 14 & 14 & 756 & 0 & 0\\
11 & 5082  & 48  & 32 & 3410 & 16  & 1672\\
13 & 22308 & 146 & 7  & 936  & 139 & 21372\\
16 & 128128 & 536 & 0 & 0 & 536 & 128128\\
17 & 210392 & 777 & 0 & 0 & 777 & 210392\\
19 & 515508 & 1514 & 0 & 0 & 1514 & 515508\\
23 & 2321781 & 4596 & 0 & 0 & 4596 & 2321781\\
25 & 4427500 & 7391 & 0 & 0 & 7391 & 4427500\\
27 & 7992270 & 11432 & 0 & 0 & 11432 & 7992270\\
29 & 13775580 & 16978 & 0 & 0 & 16978 & 13775580\\
31 & 22824711 & 24562 & 0 & 0 & 24562 & 22824711\\
32 & 28998144 & 29232 & 0 & 0 & 29232 & 28998144\\
37 & 86017008 & 64604 & 0 & 0 & 64604 & 86017008
\\
41 & 184351908 & 112439 & 0 & 0 & 112439 & 184351908
\\
43 & 262147522 & 145190 & 0 & 0 & 145190 & 262147522
\\
47 & 504665931 & 233464 & 0 & 0 & 233464 & 504665931
\\
49 & 685206984 & 291377 & 0 & 0 & 291377 & 685206984
\\
53 & 1216746440 & 441540 & 0 & 0 & 441540 & 1216746440
\end{array}
\]
}
In all twenty complete classifications, \(\rho=3\) holds if and only if the
\(56\) trisecant planes cover \(\PG(3,q)\).
\end{theorem}

\begin{proof}
The enumeration is finite and exact.  For \(q=7,8,9\), every affine
representative is constructed over its finite field and its covering radius
is computed directly; this also verifies the assertion that no radius-two
class occurs.  For \(q=11,13,17\), every affine
representative is constructed in Magma and the trisecant-plane cover is
computed directly; for \(q=11,13\) the covering radius computation was also
run as an independent check.  For \(q=19,23,29,31,37\), where no hole-free
class occurs, the scripts record an explicit uncovered projective syndrome
point for every affine representative.  The new scans for
\(q=16,25,27,32,41,43,47,49,53\) first test one or two structured projective
slices and record a point from a full-space fallback whenever both slices
fail.  The extension-field scans use field-element canonical representatives,
not integer encodings.  No class remains uncertified.  Proposition
~\ref{prop:criterion} then gives \(\rho=4\) for every certified
representative.  Reproducibility files are listed in
Section~\ref{sec:reproducibility}.
\end{proof}

\begin{theorem}[Complete covering-radius threshold]\label{thm:threshold}
For every prime power \(q\ge16\), every Roth--Lempel code
\(\RL_{4,\delta}(S)\) with \(|S|=6\) has covering radius \(4\).
For \(7\le q<16\), the complete radius-three class counts are those in
Theorem~\ref{thm:complete}: respectively \(2,3,14,32,7\) for
\(q=7,8,9,11,13\).
\end{theorem}

\begin{proof}
For \(q=16,17,19\), Proposition~\ref{prop:twisted-centers} and direct
evaluation of its one-parameter forbidden set leave only \(12,15,3\)
affine representatives, respectively.  They are certified by the explicit
families in Remark~\ref{rem:boundary-families}; the thirty certificates are
listed in the companion supplement.  Proposition~\ref{prop:twisted-centers}
proves the claim theoretically for every \(q\ge23\), which covers every
remaining prime power at least \(16\).  The final assertion is the first
five rows of Theorem~\ref{thm:complete}.
\end{proof}

\begin{corollary}[Explicit deep holes for all \(q\ge16\)]
\label{cor:explicit-deepholes}
Let \(q\ge16\), let \(C=\RL_{4,\delta}(S)\) with \(|S|=6\), and let \(u\)
be a word with nonzero syndrome \(P=Hu^{\mathsf T}\), using the explicit
matrix \(H\) in Proposition~\ref{prop:duality}.  Then \(u\) is a deep hole
of \(C\) if and only if
\[
  \Phi_{S,\delta}(P)\ne0.
\]
\end{corollary}

\begin{proof}
Theorem~\ref{thm:threshold} gives \(\rho(C)=4\) for every \(q\ge16\).
Proposition~\ref{prop:explicit-product} characterizes exactly those
syndrome cosets whose distance is \(4\).
\end{proof}

The three smallest fields have particularly compact exceptional
descriptions.  Here \(\F_8=\F_2(w)\), \(w^3+w+1=0\), and
\(\F_9=\F_3(w)\), \(w^2+2w+2=0\); all displayed classes have \(\rho=3\).
\[
\begin{array}{c|c|c|c|r}
q & \hbox{type} & S & \delta & \hbox{number of classes}\\
\hline
7 & \mathrm{NMDS} & \{0,1,2,3,4,5\} & 0,4 & 2\\
8 & \mathrm{NMDS} & \{0,1,w,w^2,w^3,w^4\} & 0,1,w^2 & 3\\
9 & \mathrm{NMDS} & \{0,1,w,w^2,w^3,2\} & \F_9 & 9\\
9 & \mathrm{NMDS} & \{0,1,w,w^2,2,w^7\} & 0,w,w^2,w^3 & 4\\
9 & \mathrm{MDS} & \{0,1,w,w^2,2,w^7\} & 1 & 1
\end{array}
\]
The unique \(\F_8\) class with \(\rho=4\) has the same displayed set
\(S=\{0,1,w,w^2,w^3,w^4\}\) and \(\delta=w^5\).

The multiplicity \(r_S(\delta)\) from Corollary~\ref{cor:duality-mds}
refines the MDS/NMDS labels.  Its complete low-field radius spectrum is:
\[
\begin{array}{c|c|r|r}
q&r_S(\delta)&\rho&\hbox{affine classes}\\ \hline
7&2&3&1\\
7&3&3&1\\ \hline
8&2&3&3\\
8&4&4&1\\ \hline
9&0&3&1\\
9&2&3&8\\
9&3&3&5\\ \hline
11&0&3&1\\
11&0&4&1\\
11&1&3&14\\
11&2&3&16\\
11&2&4&9\\
11&3&3&1\\
11&3&4&6\\ \hline
13&0&3&3\\
13&0&4&7\\
13&1&3&4\\
13&1&4&56\\
13&2&4&64\\
13&3&4&12
\end{array}
\]
Thus \(r_S(\delta)=0\) distinguishes MDS from NMDS exactly, but neither
the code type nor the number of dependent four-column sets determines the
covering radius.  For instance, over \(\F_{11}\) both radius values occur
when \(r_S(\delta)=0,2,\) or \(3\); over \(\F_{13}\), all radius-three
classes have \(r_S(\delta)\le1\).
In contrast, Proposition~\ref{prop:second-shell} shows that this same
additive statistic determines the entire second syndrome shell exactly.

Proposition~\ref{prop:deep-census} also turns the exceptional classification
into an exact deep-hole census.  The complete spectrum for the classes with
\(\rho=3\) is as follows; ``DH cosets'' denotes
\((q-1)m_3(C)\), and every such coset contains \(q^4\) deep-hole words.
{\small
\[
\begin{array}{c|c|r|r|r}
q & \hbox{type} & \hbox{DH cosets} & \hbox{affine classes} & \hbox{raw cases}\\
\hline
7  & \mathrm{NMDS} & 1380  & 1  & 7\\
7  & \mathrm{NMDS} & 1398  & 1  & 42\\
8  & \mathrm{NMDS} & 2709  & 3  & 168\\
9  & \mathrm{MDS}  & 4704  & 1  & 24\\
9  & \mathrm{NMDS} & 4752  & 8  & 516\\
9  & \mathrm{NMDS} & 4776  & 5  & 216\\
11 & \mathrm{MDS}  & 11760 & 1  & 55\\
11 & \mathrm{NMDS} & 11790 & 14 & 1540\\
11 & \mathrm{NMDS} & 11820 & 16 & 1705\\
11 & \mathrm{NMDS} & 11850 & 1  & 110\\
13 & \mathrm{MDS}  & 24432 & 3  & 312\\
13 & \mathrm{NMDS} & 24468 & 4  & 624
\end{array}
\]
}
For the radius-four classes, the hole distributions below constitute an
equally complete deep-hole census through the formula
\[
  |\mathrm{DH}(C)|=q^4(q-1)\cdot\mathrm{holes}(C).
\]
For example, the unique radius-four class over \(\F_8\) has one hole and
therefore \(8^4\cdot7=28672\) deep-hole words.

For \(q=11\) and \(q=13\), where hole-free classes occur, the distribution by
code type and by the number of uncovered projective syndrome points is:
\[
\begin{array}{c|c|c|r|r}
q & \hbox{type} & \hbox{holes} & \hbox{affine classes} & \hbox{raw cases}\\
\hline
11 & \mathrm{MDS}  & 0 & 1  & 55\\
11 & \mathrm{NMDS} & 0 & 31 & 3355\\
11 & \mathrm{MDS}  & 1 & 1  & 22\\
11 & \mathrm{NMDS} & 1 & 13 & 1430\\
11 & \mathrm{NMDS} & 2 & 2  & 220\\
\hline
13 & \mathrm{MDS}  & 0  & 3  & 312\\
13 & \mathrm{NMDS} & 0  & 4  & 624\\
13 & \mathrm{MDS}  & 1  & 4  & 624\\
13 & \mathrm{MDS}  & 2  & 1  & 156\\
13 & \mathrm{MDS}  & 3  & 2  & 312\\
13 & \mathrm{NMDS} & 1  & 11 & 1716\\
13 & \mathrm{NMDS} & 2  & 34 & 5304\\
13 & \mathrm{NMDS} & 3  & 22 & 3328\\
13 & \mathrm{NMDS} & 4  & 27 & 4082\\
13 & \mathrm{NMDS} & 5  & 15 & 2340\\
13 & \mathrm{NMDS} & 6  & 10 & 1560\\
13 & \mathrm{NMDS} & 7  & 8  & 1248\\
13 & \mathrm{NMDS} & 8  & 2  & 234\\
13 & \mathrm{NMDS} & 9  & 2  & 312\\
13 & \mathrm{NMDS} & 12 & 1  & 156
\end{array}
\]
Here holes means the number of projective syndrome points not covered by the
\(56\) trisecant planes.  Thus the rows with holes \(0\) are precisely the
classes with \(\rho=3\), and every remaining row has \(\rho=4\).

For \(q=17,19,23,29,31,37\) the classification by MDS/NMDS type is purely negative:
\[
\begin{array}{c|c|c|r|r}
q & \hbox{type} & \rho & \hbox{affine classes} & \hbox{raw cases}\\
\hline
17 & \mathrm{MDS}  & 4 & 139 & 37128\\
17 & \mathrm{NMDS} & 4 & 638 & 173264\\
19 & \mathrm{MDS}  & 4 & 350 & 119434\\
19 & \mathrm{NMDS} & 4 & 1164 & 396074\\
23 & \mathrm{MDS}  & 4 & 1488 & 752268\\
23 & \mathrm{NMDS} & 4 & 3108 & 1569513\\
29 & \mathrm{MDS}  & 4 & 7326 & 5947514\\
29 & \mathrm{NMDS} & 4 & 9652 & 7828066\\
31 & \mathrm{MDS}  & 4 & 11316 & 10518771\\
31 & \mathrm{NMDS} & 4 & 13246 & 12305940\\
37 & \mathrm{MDS}  & 4 & 34472 & 45890064\\
37 & \mathrm{NMDS} & 4 & 30132 & 40126944
\end{array}
\]
No affine class is hole-free for \(q=17,19,23,29,31,37\).  For \(q=17\), where full
hole counts were recorded, the number of uncovered projective syndrome points
ranges from \(28\) to \(117\) across the \(777\) classes.  For \(q=19,23,29,31,37\),
the original computation records an explicit first uncovered syndrome point
for every affine class.

The following earlier certificate computation searches first in the
structured projective slice \(P=[1,u,v,v]\) and invokes full-space search only
when that slice contains no hole.  Its final column is zero in every row.
\[
\begin{array}{c|r|r|r|r}
q & \hbox{affine classes} & P=[1,u,v,v]\hbox{ cert.}
  & \hbox{fallback cert.} & \hbox{uncertified}\\
\hline
17 & 777    & 695    & 82   & 0\\
19 & 1514   & 1415   & 99   & 0\\
23 & 4596   & 4396   & 200  & 0\\
29 & 16978  & 16164  & 814  & 0\\
31 & 24562  & 23434  & 1128 & 0\\
37 & 64604  & 61632  & 2972 & 0\\
41 & 112439 & 107526 & 4913 & 0
\end{array}
\]

A strengthened two-slice scan first searches \(P=[1,u,v,v]\), then
\(P=[1,u,v,u]\), and invokes full-space search only if both slices fail.
It handles both prime and non-prime fields; the final column again vanishes.
\[
\begin{array}{c|r|r|r|r|r}
q & \hbox{affine classes} & [1,u,v,v]\hbox{ cert.}
  & [1,u,v,u]\hbox{ cert.} & \hbox{fallback cert.} & \hbox{uncertified}\\
\hline
16 & 536    & 435    & 74    & 27  & 0\\
17 & 777    & 695    & 71    & 11  & 0\\
25 & 7391   & 6787   & 568   & 36  & 0\\
27 & 11432  & 10563  & 822   & 47  & 0\\
32 & 29232  & 27259  & 1866  & 107 & 0\\
43 & 145190 & 138770 & 6128  & 292 & 0\\
47 & 233464 & 224118 & 8966  & 380 & 0\\
49 & 291377 & 279078 & 11914 & 385 & 0\\
53 & 441540 & 425601 & 15368 & 571 & 0
\end{array}
\]

The \(\F_{11}\) classes with \(\rho=3\) can be summarized compactly by affine
representative families.
\[
\begin{array}{c|c|c}
\hbox{type} & S & \delta\\
\hline
\mathrm{MDS} & \{0,1,2,3,4,10\} & 10\\
\mathrm{NMDS} & \{0,1,2,3,4,10\} & 0,1\\
\mathrm{NMDS} & \{0,1,2,3,4,6\} & 1,2,3,4,5,10\\
\mathrm{NMDS} & \{0,1,2,3,4,7\} & 1,2,3,4,5,6,7,8,9,10\\
\mathrm{NMDS} & \{0,1,2,3,5,6\} & 1,2,4,5,6,7\\
\mathrm{NMDS} & \{0,1,2,3,5,9\} & 0,1,2,3,4,10\\
\mathrm{NMDS} & \{0,1,2,4,5,7\} & 2
\end{array}
\]
These rows account for \(32\) affine classes.

\begin{corollary}[The seven \(\F_{13}\) classes]\label{cor:q13-seven}
Over \(\F_{13}\), exactly the following seven affine classes have covering
radius \(3\).
\[
\begin{array}{c|c|c|c}
\hbox{type} & S & \delta & \hbox{orbit size}\\
\hline
\mathrm{MDS}  & \{0,1,2,10,11,12\} & 5  & 78\\
\mathrm{MDS}  & \{0,1,2,3,10,12\}  & 7  & 156\\
\mathrm{MDS}  & \{0,1,2,3,5,11\}   & 11 & 78\\
\mathrm{NMDS} & \{0,1,2,3,4,8\}    & 10 & 156\\
\mathrm{NMDS} & \{0,1,2,3,5,7\}    & 0  & 156\\
\mathrm{NMDS} & \{0,1,2,3,5,7\}    & 5  & 156\\
\mathrm{NMDS} & \{0,1,2,3,5,9\}    & 11 & 156
\end{array}
\]
All remaining \(139\) affine classes have covering radius \(4\).
\end{corollary}

\section{Discussion}

The classification separates two related but distinct effects.  The MDS or
NMDS status of a Roth--Lempel code is governed by additive constraints on
subsets of \(S\).  The covering-radius drop from \(4\) to \(3\), however, is a
global covering property of the eight parity-check points in \(\PG(3,q)\).
In the language of Proposition~\ref{prop:projection}, the classes with
\(\rho=4\) admit at least one projection center producing a plane \(8\)-arc;
the exceptional classes with \(\rho=3\) admit no such center.

This viewpoint suggests a finite-geometric version of the general problem:
classify eight-point configurations in \(\PG(3,q)\) arising from
\(\RL_{4,\delta}(S)\) whose \(56\) trisecant planes cover the ambient space.
The complete results above show that this cover property occurs throughout
\(q=7,9\), in three of four classes at \(q=8\), is common at \(q=11\), rare
at \(q=13\), and absent in every complete affine classification over
prime-power fields \(16\le q\le53\); the seven exceptional classes over
\(\F_{13}\) are listed in Corollary~\ref{cor:q13-seven}.
Corollary~\ref{cor:large-q} makes this rigorous for arbitrary \([8,4]_q\)
codes once \(q\ge56\), while Proposition~\ref{prop:large-q-refined} lowers
the relevant thresholds to \(q\ge54\) for MDS codes and \(q\ge53\) for NMDS
codes.  More sharply for the Roth--Lempel family itself,
Proposition~\ref{prop:twisted-centers} replaces this tail argument by
explicit deep-hole centers for every \(q\ge23\), and for every NMDS member
already when \(q\ge19\).  Consequently Theorem~\ref{thm:threshold} is a complete statement for
the entire length-eight family: a covering-radius drop can occur only below
\(q=16\), and it occurs exactly in the classes enumerated over
\(q=7,8,9,11,13\).
Proposition~\ref{prop:explicit-product} and
Corollary~\ref{cor:explicit-deepholes} add a complementary infinite
conclusion: once \(q\ge16\), the deep holes are not merely known to exist;
their syndrome cosets are given by one explicit nonvanishing condition in
four variables.
Theorem~\ref{thm:distance-stratification} removes the remaining descriptive
gap below this threshold: in the exceptional radius-three classes, the deep
holes are given explicitly by exclusion from the \(28\) displayed secant
lines.  Finally, the \(r_S(\delta)\)-spectrum shows why MDS/NMDS status alone
cannot replace the trisecant-cover condition, even though
Proposition~\ref{prop:second-shell} shows that it completely controls the
second shell.  Proposition~\ref{prop:affine-covariance} also ensures that
the new twisted-cubic certificates are intrinsic to affine classes rather
than artifacts of the chosen representatives.
Corollary~\ref{cor:closed-shells} further replaces the radius-three
deep-hole census by a closed formula depending only on \(q\) and
\(r_S(\delta)\).
From the matroidal viewpoint on linear codes, where the Tutte polynomial
records weight-enumerator information \cite{Greene1976}, the exact formula
for \(m_2(C)\) suggests that these length-eight Roth--Lempel column matroids
form a particularly rigid additive subfamily.

For reference, exact evaluation of the displayed certificate families on
the three boundary fields and the first automatic field gives
\[
\begin{array}{c|r|r|r|r}
q&\hbox{affine classes}&V(x)\hbox{ certified}&
  V(x)+\mu F\hbox{ additional}&Q(t,u)\hbox{ residual}\\ \hline
16&536&524&11&1\\
17&777&762&15&0\\
19&1514&1511&3&0\\
23&4596&4596&0&0
\end{array}
\]
Here \(Q(t,u)=(0,1,t,u)\).  The sole final residual before this third
family is the MDS class over \(\F_{16}\) with
\[
 S=\{1,w^{12},w,w^{14},0,w^8\},\qquad
 \delta=w,\qquad Q(0,w^2)=(0,1,0,w^2).
\]
These two-parameter families are used only as explicit finite certificates;
no stronger uniform threshold theorem is asserted for them.
The complete thirty-row list is supplied in
\path{Roth_Lempel_boundary_certificates_codex_supplement.tex}.

\section{Conclusion}

For the length-eight Roth--Lempel family
\(\RL_{4,\delta}(S)\), \(|S|=6\), the covering-radius problem is completely
settled.  Every member is MDS or NMDS, the exact threshold is
\(\rho(C)=4\) for all prime powers \(q\ge16\), and the only radius-three
classes occur over \(q=7,8,9,11,13\), where they are all explicitly
classified up to affine equivalence.  Moreover, the deep-hole syndromes are
described by an explicit \(56\)-factor condition in the maximal-radius
range, while the exceptional low-field cases are covered by the explicit
syndrome stratification and shell formulas.  Thus the paper gives a complete
answer for this \([8,4]\) Roth--Lempel family and provides certified seed
codes for possible deep-hole-based extension constructions; the broader
problem for arbitrary Roth--Lempel parameters \((k,n)\) remains separate.

\section{Reproducibility}\label{sec:reproducibility}

All finite computations reported in this paper were performed in Magma
\cite{Magma1997}, using exact finite-field arithmetic.  Parameter pairs are
reduced under the affine action
\[
 (S,\delta)\longmapsto (uS+v,u\delta+3v),
 \qquad u\in\F_q^{\ast},\quad v\in\F_q.
\]
No randomized search is used for a classification or for a certificate
entering a theorem.

\subsection*{Proof-critical computations}
Only two finite verification tasks enter the main results.  First, for
\(q=7,8,9,11,13\), exhaustive enumeration of affine representatives gives
the exceptional-field covering radii, distance shells, and deep-hole spectra
used in the complete small-field classification.  Second, in the three
boundary fields \(q=16,17,19\), the theoretical criterion of
Proposition~\ref{prop:twisted-centers} resolves every representative except
\(12,15,3\), respectively.  The remaining thirty representatives are
settled by explicit centers from the families \(W(x,\mu)\) and \(Q(t,u)\)
in Remark~\ref{rem:boundary-families}; the certificates themselves are
listed in the companion supplement
\path{Roth_Lempel_boundary_certificates_codex_supplement.tex}.
For every \(q\geq 23\), Theorem~\ref{thm:threshold} is theoretical and does
not depend on an exhaustive machine scan.

\subsection*{Verification summary}
\begin{center}
\small
\begin{tabular}{>{\raggedright\arraybackslash}p{0.25\linewidth}
                >{\raggedright\arraybackslash}p{0.18\linewidth}
                >{\raggedright\arraybackslash}p{0.47\linewidth}}
\toprule
Task & Fields & Certified output\\
\midrule
Exceptional classification &
\(7,8,9,11,13\) &
All affine classes; covering radius; complete distance shells and
deep-hole spectra.\\
Boundary certification &
\(16,17,19\) &
All residual cases after the theoretical test; thirty explicit
radius-four witnesses.\\
Formula validation &
\(7,\ldots,19\) &
Product criterion, shell formula, and affine-covariance checks.\\
Independent cross-checks &
\(\begin{gathered}23,25,27,29,\\
31,32,37,41,\\
43,47,49,53\end{gathered}\) &
Additional exact scans consistent with the theorem; not proof inputs.\\
\bottomrule
\end{tabular}
\end{center}

The accompanying program-and-log manifest,
\path{Roth_Lempel_reproducibility_manifest_codex_20260526.md}, records the
Magma files, exact output logs, extracted tables, and the logical role of
each computation.  This keeps the article-level argument separate from
archival computational material while providing a direct audit trail for
every finite assertion used above.

\subsection*{Data and code availability}
The Magma programs, output logs, boundary certificates, extracted data
tables, manuscript source, and reproducibility manifest are publicly
available at
\begin{center}
\url{https://github.com/matheorist/roth-lempel-covering-radius}.
\end{center}

\section{Complete \texorpdfstring{\(q=11\)}{q=11} affine classification}

This section records the complete affine classification for
\(\RL_{4,\delta}(S)\) over \(\F_{11}\), with \(|S|=6\).  There are \(48\)
affine classes representing \(5082\) raw parameter pairs \((S,\delta)\).
Exactly \(32\) affine classes, representing \(3410\) raw pairs, have
\(\rho=3\); the remaining \(16\) affine classes, representing \(1672\) raw
pairs, have \(\rho=4\).

\[
\begin{array}{c|c|c|r|r}
\hbox{type} & \rho & \hbox{holes} & \hbox{affine classes} & \hbox{raw cases}\\
\hline
\mathrm{MDS}  & 3 & 0 & 1  & 55\\
\mathrm{NMDS} & 3 & 0 & 31 & 3355\\
\mathrm{MDS}  & 4 & 1 & 1  & 22\\
\mathrm{NMDS} & 4 & 1 & 13 & 1430\\
\mathrm{NMDS} & 4 & 2 & 2  & 220
\end{array}
\]

Here holes denotes the number of projective syndrome points not covered by
the \(56\) trisecant planes.  Thus the hole-free rows are exactly the
\(\rho=3\) rows.

\begin{longtable}{r r c c l r r r}
\caption{All affine classes for \(q=11,k=4,n=6\).}\\
\toprule
idx & orbit & type & \(\rho\) & \(S\) & \(\delta\) & min & holes\\
\midrule
\endfirsthead
\toprule
idx & orbit & type & \(\rho\) & \(S\) & \(\delta\) & min & holes\\
\midrule
\endhead
1 & 110 & NMDS & 3 & \(\{0,1,2,3,5,6\}\) & 7 & 1 & 0\\
2 & 110 & NMDS & 3 & \(\{0,1,2,3,5,6\}\) & 5 & 1 & 0\\
3 & 110 & NMDS & 4 & \(\{0,1,2,3,5,6\}\) & 3 & 0 & 1\\
4 & 110 & NMDS & 3 & \(\{0,1,2,3,5,6\}\) & 1 & 1 & 0\\
5 & 110 & NMDS & 4 & \(\{0,1,2,3,5,6\}\) & 10 & 0 & 1\\
6 & 110 & NMDS & 4 & \(\{0,1,2,3,5,6\}\) & 8 & 0 & 1\\
7 & 110 & NMDS & 3 & \(\{0,1,2,3,5,6\}\) & 6 & 1 & 0\\
8 & 110 & NMDS & 3 & \(\{0,1,2,3,5,6\}\) & 4 & 1 & 0\\
9 & 110 & NMDS & 3 & \(\{0,1,2,3,5,6\}\) & 2 & 1 & 0\\
10 & 110 & NMDS & 4 & \(\{0,1,2,3,5,6\}\) & 0 & 0 & 1\\
11 & 110 & NMDS & 4 & \(\{0,1,2,3,5,6\}\) & 9 & 0 & 1\\
12 & 110 & NMDS & 3 & \(\{0,1,2,3,4,7\}\) & 9 & 1 & 0\\
13 & 110 & NMDS & 3 & \(\{0,1,2,3,4,7\}\) & 1 & 1 & 0\\
14 & 110 & NMDS & 3 & \(\{0,1,2,3,4,7\}\) & 4 & 1 & 0\\
15 & 110 & NMDS & 3 & \(\{0,1,2,3,4,7\}\) & 7 & 1 & 0\\
16 & 110 & NMDS & 3 & \(\{0,1,2,3,4,7\}\) & 10 & 1 & 0\\
17 & 110 & NMDS & 3 & \(\{0,1,2,3,4,7\}\) & 2 & 1 & 0\\
18 & 110 & NMDS & 3 & \(\{0,1,2,3,4,7\}\) & 5 & 1 & 0\\
19 & 110 & NMDS & 3 & \(\{0,1,2,3,4,7\}\) & 8 & 1 & 0\\
20 & 110 & NMDS & 4 & \(\{0,1,2,3,4,7\}\) & 0 & 0 & 1\\
21 & 110 & NMDS & 3 & \(\{0,1,2,3,4,7\}\) & 3 & 1 & 0\\
22 & 110 & NMDS & 3 & \(\{0,1,2,3,4,7\}\) & 6 & 1 & 0\\
23 & 110 & NMDS & 4 & \(\{0,1,2,3,4,6\}\) & 6 & 0 & 1\\
24 & 110 & NMDS & 3 & \(\{0,1,2,3,4,6\}\) & 4 & 1 & 0\\
25 & 110 & NMDS & 3 & \(\{0,1,2,3,4,6\}\) & 2 & 1 & 0\\
26 & 110 & NMDS & 4 & \(\{0,1,2,3,4,6\}\) & 0 & 0 & 1\\
27 & 110 & NMDS & 4 & \(\{0,1,2,3,4,6\}\) & 9 & 0 & 1\\
28 & 110 & NMDS & 4 & \(\{0,1,2,3,4,6\}\) & 7 & 0 & 1\\
29 & 110 & NMDS & 3 & \(\{0,1,2,3,4,6\}\) & 5 & 1 & 0\\
30 & 110 & NMDS & 3 & \(\{0,1,2,3,4,6\}\) & 3 & 1 & 0\\
31 & 110 & NMDS & 3 & \(\{0,1,2,3,4,6\}\) & 1 & 1 & 0\\
32 & 110 & NMDS & 3 & \(\{0,1,2,3,4,6\}\) & 10 & 1 & 0\\
33 & 110 & NMDS & 4 & \(\{0,1,2,3,4,6\}\) & 8 & 0 & 1\\
34 & 110 & NMDS & 4 & \(\{0,1,2,3,4,10\}\) & 4 & 0 & 2\\
35 & 110 & NMDS & 3 & \(\{0,1,2,3,4,10\}\) & 0 & 1 & 0\\
36 & 110 & NMDS & 4 & \(\{0,1,2,3,4,10\}\) & 2 & 0 & 1\\
37 & 110 & NMDS & 4 & \(\{0,1,2,3,4,10\}\) & 3 & 0 & 2\\
38 & 55 & MDS & 3 & \(\{0,1,2,3,4,10\}\) & 10 & 1 & 0\\
39 & 110 & NMDS & 3 & \(\{0,1,2,3,4,10\}\) & 1 & 1 & 0\\
40 & 110 & NMDS & 3 & \(\{0,1,2,3,5,9\}\) & 2 & 1 & 0\\
41 & 110 & NMDS & 3 & \(\{0,1,2,3,5,9\}\) & 1 & 1 & 0\\
42 & 110 & NMDS & 3 & \(\{0,1,2,3,5,9\}\) & 0 & 1 & 0\\
43 & 55 & NMDS & 3 & \(\{0,1,2,3,5,9\}\) & 10 & 1 & 0\\
44 & 110 & NMDS & 3 & \(\{0,1,2,3,5,9\}\) & 3 & 1 & 0\\
45 & 110 & NMDS & 3 & \(\{0,1,2,3,5,9\}\) & 4 & 1 & 0\\
46 & 110 & NMDS & 4 & \(\{0,1,2,4,5,7\}\) & 0 & 0 & 1\\
47 & 110 & NMDS & 3 & \(\{0,1,2,4,5,7\}\) & 2 & 1 & 0\\
48 & 22 & MDS & 4 & \(\{0,1,2,4,5,7\}\) & 4 & 0 & 1\\
\bottomrule
\end{longtable}

\noindent\textbf{Companion appendices.}
Complete affine-class tables for \(q=13,17,19,23,29,31,37\) are supplied as
Appendices B--H in the companion appendix volume:
\begin{center}
\footnotesize\path{Roth_Lempel_trisecant_covers_DM_codex_appendices_BH_q13_q37.pdf}.
\end{center}
For \(q=19,23,29,31,37\), each listed class includes an explicit first-hole
certificate.  The full class-level results for \(q=7,8,9\) are contained in
the exact small-prime-power log.  Exact certificate logs for the later fields
\(q=41,43,47,49,53\) and the extension fields \(q=16,25,27,32\) are listed
in Section~\ref{sec:reproducibility} and are available in the public
repository above.

\begin{thebibliography}{99}

\bibitem{Ball2015}
S. Ball,
\emph{Finite Geometry and Combinatorial Applications},
Cambridge University Press, 2015.

\bibitem{Cohen1997}
G. Cohen, I. Honkala, S. Litsyn, A. Lobstein,
\emph{Covering Codes},
North-Holland Mathematical Library, vol. 54, Elsevier, 1997.

\bibitem{DodunekovLandgev1995}
S. M. Dodunekov, I. N. Landgev,
``On near-MDS codes,''
\emph{J. Geom.}, vol. 54, no. 1--2, pp. 30--43, 1995,
doi:10.1007/BF01222850.

\bibitem{Greene1976}
C. Greene,
``Weight enumeration and the geometry of linear codes,''
\emph{Stud. Appl. Math.}, vol. 55, no. 2, pp. 119--128, 1976,
doi:10.1002/sapm1976552119.

\bibitem{HanFan2023}
D. Han, C. Fan,
``Roth--Lempel NMDS codes of non-elliptic-curve type,''
\emph{IEEE Trans. Inf. Theory}, vol. 69, no. 9, pp. 5670--5675, 2023,
doi:10.1109/TIT.2023.3272384.

\bibitem{Huffman2003}
W. C. Huffman, V. Pless,
\emph{Fundamentals of Error-Correcting Codes},
Cambridge University Press, 2003.

\bibitem{LiLuLingLam2026}
Y. Li, Z. Lu, S. Ling, K.-Y. Lam,
``A framework for constructing non-GRS MDS-NMDS codes from deep holes and its application,''
arXiv:2605.12133v1, May 2026. \url{https://arxiv.org/abs/2605.12133}.

\bibitem{LiZhuSun2025}
Y. Li, S. Zhu, Z. Sun,
``Covering radii and deep holes of two classes of extended twisted GRS codes
and their applications,''
\emph{IEEE Trans. Inf. Theory}, vol. 71, no. 5, pp. 3516--3530, 2025,
doi:10.1109/TIT.2025.3541799.

\bibitem{LiangLiao2026}
Z. Liang, Q. Liao,
``Two classes of NMDS codes from Roth--Lempel codes,''
\emph{Finite Fields Appl.}, vol. 111, Art. no. 102779, 2026,
doi:10.1016/j.ffa.2025.102779.

\bibitem{Magma1997}
W. Bosma, J. Cannon, C. Playoust,
``The Magma algebra system I: The user language,''
\emph{J. Symb. Comput.}, vol. 24, no. 3--4, pp. 235--265, 1997,
doi:10.1006/jsco.1996.0125.

\bibitem{MWS1977}
F. J. MacWilliams, N. J. A. Sloane,
\emph{The Theory of Error-Correcting Codes},
North-Holland, Amsterdam, 1977.

\bibitem{RothLempel1989}
R. M. Roth, A. Lempel,
``A construction of non-Reed-Solomon type MDS codes,''
\emph{IEEE Trans. Inf. Theory}, vol. 35, no. 3, pp. 655--657, 1989,
doi:10.1109/18.30988.

\bibitem{XuZhou2026}
H. Xu, H. Zhou,
``Analysis of Roth--Lempel codes,''
\emph{IEEE Trans. Inf. Theory}, vol. 72, no. 1, pp. 246--252, 2026,
doi:10.1109/TIT.2025.3626824.

\end{thebibliography}

\end{document}
